Kiekvienas priedas pradedamas makrokomanda \verb|\priedas{Pavadinimas}|, kuri sukuria antraštę „N priedas. Pavadinimas“ ir įtraukia ją į turinį. Priedų lentelės ir paveikslai numeruojami „P1“, „P2“ ir t.\,t.

\priedas{Pirmojo priedo pavadinimas}

\lipsum[1]

\begin{table}[!ht]
\caption{Pavyzdinė priedo lentelė}
\label{tab:priedas-pavyzdys}
\centering
\begin{tabular}{|p{6cm}|r|} \toprule
\textbf{Rodiklis} & \textbf{Reikšmė} \\ \midrule
Rodiklis 1 & 100{,}0 \\ \hline
Rodiklis 2 & 42{,}0 \\ \bottomrule
\end{tabular}
\end{table}

\priedas{Antrojo priedo pavadinimas}

\lipsum[2]
